\documentclass[10pt,a4paper,openany,twoside]{book}

\usepackage{layout_cours}

\begin{document}

\titre{Equation de degré un et équation produit nul}
\classe{Seconde}
\thispagestyle{premierepage}
\enteteperso

\pagenumbering{Alph}
\section*{Bordas 2019}
\begin{exercicesnofoot}
\hspace{-1.5cm}
\includegraphics[width=0.5\textwidth]{../../../Seconde/Bordas_Indice_2nde_2019/Bordas_Indice_2nde_2019_p_104_exo_115_116_117_118.jpg}

\hspace{-1.5cm}
\includegraphics[width=0.5\textwidth]{../../../Seconde/Bordas_Indice_2nde_2019/Bordas_Indice_2nde_2019_p_104_exo_119.jpg}

\hspace{-1.5cm}
\includegraphics[width=0.5\textwidth]{../../../Seconde/Bordas_Indice_2nde_2019/Bordas_Indice_2nde_2019_p_104_exo_120_121.jpg}

\hspace{-1.5cm}
\includegraphics[width=0.5\textwidth]{../../../Seconde/Bordas_Indice_2nde_2019/Bordas_Indice_2nde_2019_p_104_exo_123_124_125_126.jpg}

\hspace{-1.5cm}
\includegraphics[width=0.5\textwidth]{../../../Seconde/Bordas_Indice_2nde_2019/Bordas_Indice_2nde_2019_p_113_exo_197.jpg}

\end{exercicesnofoot}

\newpage
\pagenumbering{arabic}
\thispagestyle{premierepage}
\enteteperso

\pagenumbering{Alph}
\section*{Bordas 2019}
\begin{exercicesnofoot}
\hspace{-1.5cm}
\includegraphics[width=0.5\textwidth]{../../../Seconde/Bordas_Indice_2nde_2019/Bordas_Indice_2nde_2019_p_104_exo_115_116_117_118.jpg}

\hspace{-1.5cm}
\includegraphics[width=0.5\textwidth]{../../../Seconde/Bordas_Indice_2nde_2019/Bordas_Indice_2nde_2019_p_104_exo_119.jpg}

\hspace{-1.5cm}
\includegraphics[width=0.5\textwidth]{../../../Seconde/Bordas_Indice_2nde_2019/Bordas_Indice_2nde_2019_p_104_exo_120_121.jpg}

\hspace{-1.5cm}
\includegraphics[width=0.5\textwidth]{../../../Seconde/Bordas_Indice_2nde_2019/Bordas_Indice_2nde_2019_p_104_exo_123_124_125_126.jpg}

\hspace{-1.5cm}
\includegraphics[width=0.5\textwidth]{../../../Seconde/Bordas_Indice_2nde_2019/Bordas_Indice_2nde_2019_p_113_exo_197.jpg}

\end{exercicesnofoot}

\newpage
\pagenumbering{arabic}
\section*{Bordas 2019}

    \sectionexo{Equation du premier degré et équation quotient nul}
    \begin{bookexos}{115, 116 p.104}\end{bookexos}
    \begin{bookexos}{124 p.104}\end{bookexos}
    \sectionexo{Equation produit nul}
    \begin{bookexos}{117-121 p.104}\end{bookexos}
    % \begin{bookexos}{125-126 p.105}\end{bookexos}
    % \begin{bookexo}{129 p.105}\end{bookexo}
    \sectionexo{Vers les équations du second degré}
    \begin{bookexos}{122-123 p.104}\end{bookexos}

    \sectionexo{Parcours d'approfondissement}
    \begin{bookexo}{197 p.113}\end{bookexo}

\newpage
\section*{Correction}
\setcounter{sectionexo}{0}
    \sectionexo{Equation du premier degré et équation quotient nul}
    \begin{bookexo}{115 p.104}
        \begin{enumerate}[label=\alph*.]
            \item 
            $S=\{\frac{-1}{5}\}$
            \item 
            $S=\{2\}$
        \end{enumerate}
    \end{bookexo}

    \begin{bookexo}{116 p.104}
        \begin{enumerate}[label=\alph*.]
            \item 
            $S=\{\frac{-23}{14}\}$
            \item 
            $S=\{\frac{-2}{21}\}$
        \end{enumerate}
    \end{bookexo}
    \begin{bookexos}{124 p.104}
             \begin{enumerate}[label=\alph*.]
            \item 
            $S=\{3\}$
            \item 
            $S=\{2\}$
            \item 
            $S=\{\frac{1}{9}\}$
        \end{enumerate}
    \end{bookexos}
    \sectionexo{Equation produit nul}
    \begin{bookexo}{117 p.104}
             \begin{enumerate}[label=\alph*.]
            \item 
            $S=\{-3;\quad 2\}$
            \item 
            $S=\{-\frac{1}{5};\quad \frac{3}{2}\}$
        \end{enumerate}
    \end{bookexo}
    \begin{bookexo}{118 p.104}
             \begin{enumerate}[label=\alph*.]
            \item 
            $S=\{-3;\quad 0\}$
            \item 
            $S=\{-14;\quad 1\}$
        \end{enumerate}
    \end{bookexo}
    \begin{bookexo}{119 p.104}
             \begin{enumerate}
            \item On calcule $(3x-5)(2-x)=3x^2+11x-10=A(x)$
            \item 
            $S=\{2;\quad \frac{5}{3}\}$
        \end{enumerate}
    \end{bookexo}
  \begin{bookexo}{120 p.104}
             \begin{enumerate}[label=\alph*.]
            \item 
            $S=\{0;\quad 5\}$
            \item 
            $S=\{0;\quad \frac{1}{3}\}$
            \item 
            $S=\{0;\quad 7\}$
        \end{enumerate}
    \end{bookexo}
      \begin{bookexo}{121 p.104}
             \begin{enumerate}[label=\alph*.]
            \item 
            $S=\{-4;\quad \frac{1}{2}\}$
            \item 
            $S=\{\frac{-2}{3};\quad 1\}$
        \end{enumerate}
    \end{bookexo}
    % \begin{bookexos}{125-126 p.105}\end{bookexos}
    % \begin{bookexo}{129 p.105}\end{bookexo}
    \sectionexo{Vers les équations du second degré}
    \begin{bookexo}{122 p.104}
            $S=\{-1;\quad -7\}$
    \end{bookexo}
      \begin{bookexo}{123 p.104}
             \begin{enumerate}[label=\alph*.]
            \item 
            $S=\{-7;\quad 1\}$
            \item 
            $S=\{\frac{-1}{2};\quad \frac{5}{4}\}$
        \end{enumerate}
    \end{bookexo}

    \sectionexo{Parcours d'approfondissement}
    \begin{bookexo}{197 p.113}
        \begin{enumerate}
            \item $h=15.5$ cm et $g=33.8$ cm.
            \item Avec $h=16$ cm, on a $g=32.8$ cm. Soit au plus $N=9$ marches. Donc une hauteur totale de $9h=1.44$ m.
        \end{enumerate}
    \end{bookexo}
\end{document}